\BourbakiExerciseGroup

\begin{enumerate}
\def\labelenumi{\arabic{enumi})}
\item
  Extend Propositions 3, 6, 7 and 9, and the corollary to Proposition 4, to semi-topological groups (§ 1, Exercise 2) and paratopological groups (§ 1, Exercise 4).
\item
  \begin{enumerate}
  \def\labelenumii{\alph{enumii})}
  \tightlist
  \item
    Extend Propositions 1, 2 and 4 to semi-topological groups (§ 1, Exercise 2). (To prove that the closure \(\overline{\mathbf{H}}\) of a subgroup \(\mathbf{H}\) is a subgroup start by remarking that \(\overline{s\overline{\mathbf{H}}} \subset \overline{\mathbf{H}}\) for all \(s \in \mathbf{H}\) ).
  \end{enumerate}
\end{enumerate}

\begin{enumerate}
\def\labelenumi{\alph{enumi})}
\setcounter{enumi}{1}
\tightlist
\item
  Let \(\mathbf{G}\) be the group \(\mathbf{Z}^{(\mathbb{N})}\) , and let \(e_i\) ( \(i \in \mathbb{N}\) ) be the element of \(\mathbf{G}\) all of whose coordinates are zero except the \(i\) th, and whose \(i\) th coordinate is 1. For each \(n \in \mathbb{N}\) , let \(\mathbf{V}_n\) be the set of all \(z = (z_k) \in \mathbf{G}\) such that \(z_k = 0\) for \(k \leqslant n\) and \(z_k \geqslant 0\) for \(k > n\) . The \(\mathbf{V}_n\) form a fundamental system of neighbourhoods of \(0\) for a topology which makes \(\mathbf{G}\) into a Hausdorff paratopological group. Let H be the subgroup of G generated by the elements \(e_{0} + e_{n} (n \geqslant 1)\) . Show that H is discrete but not closed in G, and that \(\overline{H}\) is not a subgroup of G.
\end{enumerate}

\begin{enumerate}
\def\labelenumi{\arabic{enumi})}
\setcounter{enumi}{2}
\tightlist
\item
  Give an example of a non-Hausdorff topological group in which the centre is not closed and the closure of the centre is a non-commutative subgroup (cf.~§ 1, Exercise 1). Give an example of a semi-topological group having the same properties and in which every set consisting of a single point is closed {[}cf.~§ 1, Exercise 2 b){]}.
\end{enumerate}

¶ 4) a) In a semi-topological group G, the intersection of the neighbourhoods of e which are both open and closed is a closed subgroup H which is invariant under all automorphisms of G (remark that there can be no partition of G into two open and closed subsets, each of which meets H).

\begin{enumerate}
\def\labelenumi{\alph{enumi})}
\setcounter{enumi}{1}
\tightlist
\item
  On the set \(\mathfrak{G}\) of all subgroups of G, define a mapping \(\varphi: \mathfrak{G} \to \mathfrak{G}\) as follows: if H is any subgroup of G, then \(\varphi(H)\) is the subgroup of H which is the intersection of all the simultaneously open and closed neighbourhoods of \(e\) in H. In the set \(\mathfrak{G}\) , ordered by the relation \(\supset\) , consider the chain \(\Gamma\) of G with respect to the mapping \(\varphi\) (Set Theory, Chapter III, § 2, Exercise 6); show that the smallest subgroup of \(\Gamma\) is the connected component of \(e\) in G (remark that this smallest subgroup must be connected).
\end{enumerate}

\begin{enumerate}
\def\labelenumi{\arabic{enumi})}
\setcounter{enumi}{4}
\tightlist
\item
  A semi-topological group G is quasi-topological if, for each \(a \in G\) , the mapping \(x \to xax^{-1}\) of G into G is continuous. Every commutative semi-topological group is quasi-topological.
\end{enumerate}

\begin{enumerate}
\def\labelenumi{\alph{enumi})}
\item
  Show that the normalizer of a closed subgroup of a quasi-topological group is closed.
\item
  Let G be a quasi-topological group which is generated by each neighbourhood of e (resp. which is connected). Show that every discrete (resp. totally disconnected) normal subgroup D of G is contained in the centre of G (if \(a \in D\) , show that there is a neighbourhood V of e such that \(xax^{-1} = a\) for each \(x \in V\) ).
\end{enumerate}

\(^{*}c)\) Let \(G\) be the group of matrices \(\begin{pmatrix} 1 & 0 \\ y & x \end{pmatrix}\) , where \(x\) and \(y\) are real numbers and \(x > 0\) . If we identify \(G\) with a subset of \(R^2\) by means of the mapping \((x, y) \to \begin{pmatrix} 1 & 0 \\ y & x \end{pmatrix}\) , then the topology \(\mathfrak{G}_0\) induced on \(G\) by the topology of \(R^2\) is compatible with the group structure of \(G\) . Let \(H\) be the normal subgroup of \(G\) consisting of all matrices \(\begin{pmatrix} 1 & 0 \\ y & 1 \end{pmatrix}\) , and let \(H^*\) be the complement of the identity element in \(H\) . Define a topology \(\mathfrak{G}\) , finer than \(\mathfrak{G}_0\) , by taking as a fundamental system of neighbourhoods of \(e\) in \(G\) the intersections of neighbourhoods of \(e\) in the topology \(\mathfrak{G}_{0}\) with the complement of H* in G. With respect to this topology \(\mathfrak{G}\) , show that G is a connected semi-topological group which is not quasi-topological and which has H as a discrete subgroup not contained in the centre. *

* 6) Let G be the orthogonal group \(\mathbf{O}_{2}(\mathbf{Q})\) , identified with a subspace of the space \(\mathbf{Q}^{4}\) of all \(2 \times 2\) matrices over \(\mathbf{Q}\) . Show that the topology induced on G by the topology of \(\mathbf{Q}^{4}\) is compatible with the group structure of G. Show that the commutator subgroup of G is not closed and that its closure in G is the group \(\mathbf{SO}_{2}(\mathbf{Q})\) . *

\begin{enumerate}
\def\labelenumi{\arabic{enumi})}
\setcounter{enumi}{6}
\tightlist
\item
  \begin{enumerate}
  \def\labelenumii{\alph{enumii})}
  \tightlist
  \item
    Show that if G is a connected quasi-topological group (Exercise 5) then the commutator subgroup of G is connected. (Let \(P_k\) be the set of all products of k commutators \(x^{-1}y^{-1}xy\) ; show that \(P_k\) is a union of connected sets each of which meets \(P_{k-1}\) ).
  \end{enumerate}
\end{enumerate}

\begin{enumerate}
\def\labelenumi{\alph{enumi})}
\setcounter{enumi}{1}
\tightlist
\item
  Let G be a non-commutative infinite group whose commutator subgroup is finite (for example, the product of a non-commutative finite group and an infinite commutative group). If G carries the topology defined in Exercise 2 b) of § 1, show that G is a connected semi-topological group whose commutator subgroup is not connected.
\end{enumerate}

¶ 8) Let G be a quasi-topological group (Exercise 5).

\begin{enumerate}
\def\labelenumi{\alph{enumi})}
\item
  Extend Proposition 8 to G (show first that \(x\overline{\mathbf{K}}x^{-1} \subset \overline{\mathbf{K}}\) if \(x \in \mathbf{H}\) ). Give an example of semi-topological group which is not quasi-topological and for which Proposition 8 is not valid {[}cf.~§ 1, Exercise 2 b){]}.
\item
  Let H, K be two subgroups of G such that H⊃K and K contains the commutator subgroup of H. Show that \(\overline{K}\) contains the commutator subgroup of \(\overline{H}\) (analogous method).
\item
  Deduce from b) that if every subset of G consisting of a single point is closed, then the closure of every solvable (resp. commutative) subgroup in G is solvable (resp. commutative) (argue on the length of the derived series of the subgroup under consideration).
\end{enumerate}

\begin{enumerate}
\def\labelenumi{\arabic{enumi})}
\setcounter{enumi}{8}
\item
  Let G be a quasi-topological group, and let H be a closed normal subgroup of G which contains the commutator subgroup of G. Show that if the component K of e in H is solvable, then so is the component L of e in G {[}show by use of Exercise 7 a) that K contains the commutator subgroup of L{]}.
\item
  Let G be a quasi-topological group in which every subset consisting of a single point is closed and whose identity component C is such that G/C is finite. Show that if G/C has k elements, then the set of conjugates \(xax^{-1}\) of an element \(a \in G\) is either infinite or else has at most k elements (consider the mapping \(x \rightarrow xax^{-1}\) of G onto this set). In particular, if G is connected, the set of conjugates of every element not in the centre of G is infinite.
\item
  Let G be a group, topologized or not, which operates on a topological space X in such a way that each of the mappings \(x \rightarrow s.x\) ( \(s \in G\) ) of X into X is continuous. Extend the results of no. 4 to this situation.
\item
  Let G be a semi-topological (resp. paratopological, quasi-topological) group and let H be a normal subgroup of G.
\end{enumerate}

\begin{enumerate}
\def\labelenumi{\alph{enumi})}
\item
  Show that if G/H is endowed with the quotient topology of G by H, then G/H is a semi-topological (resp. paratopological, quasi-topological) group.
\item
  Extend the results of no. 7 to semi-topological (resp. paratopological, quasi-topological) groups.
\end{enumerate}

\begin{enumerate}
\def\labelenumi{\arabic{enumi})}
\setcounter{enumi}{12}
\item
  Let G be a semi-topological (resp. paratopological) group and let H be a dense subgroup of G. What is the topology of the homogeneous space G/H?
\item
  Let G be a semi-topological group and let H be a normal subgroup of G. Show that the semi-topological group \(G/\overline{H}\) is isomorphic to the Hausdorff group associated with G/H.
\end{enumerate}

¶ 15) Let G be a quasi-topological group in which every set consisting of a single point is closed. Show that if G is solvable then there is a composition series of G consisting of closed subgroups such that the successive quotient groups of the series are commutative {[}argue by induction on the length of the derived series of G, using Exercise 8 c){]}.

¶ 16) Let H be a subgroup of a semi-topological group G, contained in the component K of e. Show that the components of the space G/H are the images of the components of G under the canonical mapping f of G onto G/H {[}if L is a component of G/H, show that \(\overline{f}^{-1}(L)\) is connected, by reductio ad absurdum{]}. Show that K is the smallest of the closed subgroups H of G such that G/H is totally disconnected.

* 17) Let G be the additive group of mappings \(n \to f(n)\) of N into Q such that \(\lim_{n \to \infty} f(n)\) exists in R. For each \(\alpha > 0\) , let \(V_{\alpha}\) be the set of \(f \in G\) such that \(|f(n)| < \alpha\) for all \(n \in \mathbb{N}\) . Show that the \(V_{\alpha}\) satisfy axioms \((GV_{I})\) and \((GV_{II})\) and that with the topology defined by the \(V_{\alpha}\) , the group G is Hausdorff and totally disconnected. Let H be the subgroup of all \(f \in G\) such that \(\lim_{n \to \infty} f(n) = 0\) . Show that H is closed and that G/H is isomorphic to R and therefore connected. *

\begin{enumerate}
\def\labelenumi{\arabic{enumi})}
\setcounter{enumi}{17}
\item
  A continuous homomorphism \(f\) of a topological group \(G\) into a topological group \(G'\) is a strict morphism if and only if the image under \(f\) of every open set in \(G\) has at least one interior point with respect to \(f(G)\) .
\item
  Let \(G, G', G''\) be three topological groups; let \(f\) be a strict morphism of \(G\) into \(G'\) , and let \(g\) be a strict morphism of \(F'\) into \(G''\) . Show that if \(f(G)\) contains the kernel \(\overline{g}^{-1}(e'')\) of \(g\) ( \(e''\) being the identity element of \(G''\) ), then \(g \circ f\) is a strict morphism of \(G\) into \(G''\) (cf.~no. 7, Corollary to Proposition 21). Give an example in which \(f\) is injective and the preceding condition is not satisfied, and \(g \circ f\) is not a strict morphism of \(G\) into \(G''\) (cf.~no. 7, Remark following Proposition 20).
\item
  Let H, H' be two subgroups of a semi-topological group G, such that \(H' \subset H\) . Show that if \(f, f'\) are the canonical mappings of G onto G/H and \(G/H'\) respectively, then there is a unique mapping \(\varphi\) of \(G/H'\) onto G/H such that \(f = \varphi \circ f'\) ; \(\varphi\) is continuous and open. If moreover \(H'\) is open in H, then every point x of \(G/H'\) has a neighbourhood V such that the restriction of \(\varphi\) to V is a homeomorphism of V onto a neighbourhood of \(\varphi(x)\) in G/H.
\item
  Let G be a topological group and let H, K be two subgroups of G such that G = HK and \(H \cap K = \{e\}\) . Each \(x \in G\) can be written uniquely in the form \(x = f(x)g(x)\) , where \(f(x) \in H\) and \(g(x) \in K\) . The mapping \((y, z) \to yz\) of \(H \times K\) onto G is a homeomorphism if and only if one of the mappings f, g is continuous. Show that this condition is always satisfied if one of the subgroups H, K is compact and closed, and the other is closed (remark that if H and K are closed, f and g can have no cluster point other than e when x tends to e in G).
\item
  Let G be the product of a family \((\mathbf{G}_{i})_{i\in I}\) of topological groups, and let \(G'\) be a topological group such that there is a neighbourhood \(V'\) of the identity element \(e'\) in \(G'\) which contains no normal subgroup of \(G'\) other than \(\{e'\}\) . If f is a continuous homomorphism of G into \(G'\) , show that there is a subset J of I whose complement \(J'\) is finite, such that f is equal to \(e'\) on the subgroup \(\prod_{i\in I}G_{i}\times\prod_{i\in J'}\{e_{i}\}\) of G.
\end{enumerate}

\[
\left(\mathbf {G} _ {i}\right) _ {i \in \mathbf {I}}
\]

\[
\mathbf {G} _ {i}
\]

\[
\prod_ {i \in I} A _ {i},
\]

where \(\mathbf{A}_i\) is open in \(\mathbf{G}_i\) for each \(i \in I\) and where the set of \(i \in I\) such that \(\mathbf{A}_i \neq \mathbf{G}_i\) has cardinal \(< c\) , where \(c\) is a given infinite cardinal (Chapter I, § 4, Exercise 9). Show that this topology is compatible with the group structure of G. Hence construct an example of a non-discrete Hausdorff topological group in which every compact subset is finite (cf.~Chapter I, § 9, Exercise 4).

\begin{enumerate}
\def\labelenumi{\arabic{enumi})}
\setcounter{enumi}{23}
\tightlist
\item
  Let G and G' be two locally isomorphic connected groups.
\end{enumerate}

\begin{enumerate}
\def\labelenumi{\alph{enumi})}
\item
  Let \(f\) be a local isomorphism of \(\mathbf{G}\) with \(\mathbf{G}'\) , defined on a neighbourhood \(\mathbf{V}\) of the identity element of \(\mathbf{G}\) . Let \(\mathbf{H}\) be the subgroup of the product \(\mathbf{G} \times \mathbf{G}'\) generated by the set of points \((x, f(x))\) where \(x \in \mathbf{V}\) . Consider the filter base on \(\mathbf{H}\) consisting of the images of neighbourhoods of \(e\) contained in \(\mathbf{V}\) under the mapping \(x \to (x, f(x))\) . Show that this filter base is a fundamental system of neighbourhoods of the identity element of \(\mathbf{H}\) for a topology \(\mathcal{C}\) compatible with the group structure of \(\mathbf{H}\) .
\item
  Show that there exist two discrete normal subgroups K, K' contained in the centre of H, such that G is isomorphic to H/K and G' is isomorphic to H/K' {[}use Exercise 5 b){]}.
\end{enumerate}

* c) Take G and G' each to be the group T; let φ be the canonical homomorphism of R onto T = R/Z, and let f be the local isomorphism which maps each point φ(x), where \(-1/4 \leqslant x \leqslant 1/4\) , to φ(θx), θ being an irrational number such that 0 \textless{} θ \textless{} 1. Show that for this local isomorphism, the topology C defined on H is distinct from the topology induced on H by the product topology on G × G'*.

\begin{enumerate}
\def\labelenumi{\arabic{enumi})}
\setcounter{enumi}{24}
\item
  Let G be a topological group, K a normal subgroup of G, \(\varphi\) the canonical mapping \(G \rightarrow G/K\) . Let f be a continuous homomorphism of a topological group H into G, such that every element of \(f(H)\) commutes with every element of K. Then \((y, z) \rightarrow f(y)z\) is a continuous homomorphism g of \(H \times K\) into G. g is a strict morphism of \(H \times K\) into G if and only if \(\varphi \circ f\) is a strict morphism of H into G/K.
\item
  Let \((\mathbf{G}_i)_{i\in I}\) be a family of topological groups, and for each \(i\in I\) let \(\mathbf{K}_i\) be an open normal subgroup of \(\mathbf{G}_i\) . Let \(\mathfrak{B}\) denote the neighbourhood filter of the identity element in the product topological group \(\mathbf{K} = \prod_{i\in I}\mathbf{K}_i\) . Show that in the product group \(\mathbf{G} = \prod_{i\in I}\mathbf{G}_i\) , \(\mathfrak{B}\) is a base of the neighbourhood filter of \(e\) for a topology compatible with the group structure of \(\mathbf{G}\) . The group \(\mathbf{G}\) , endowed with this topology, is called the local product of the \(\mathbf{G}_i\) (relative to the \(\mathbf{K}_i\) ); \(\mathbf{K}\) is then an open normal subgroup of \(\mathbf{G}\) , and \(\mathbf{G} / \mathbf{K}\) is isomorphic to the (discrete) product of the groups \(\mathbf{G}_i / \mathbf{K}_i\) . If, for each \(i\in I\) , \(\mathbf{K}_i'\) is another open normal subgroup of \(\mathbf{G}_i\) , then the local product topologies on \(\mathbf{G}\) , defined by the families \((\mathbf{K}_i)\) and \((\mathbf{K}_i')\) , coincide if and only if \(\mathbf{K}_i' = \mathbf{K}_i\) for all but a finite number of indices.
\end{enumerate}

For each \(i \in I\) , let \(G_i'\) be the normal subgroup of \(G\) consisting of all \(x = (x_i)\) such that \(x_i = e_i\) whenever \(i \neq x\) . Then \(G_i'\) , with the topology induced by a local product topology on \(G\) , is canonically isomorphic to \(G_i\) . The closure in \(G\) of the subgroup generated by the \(G_i'\) is the normal subgroup \(G_0\) of \(G\) which consists of all \(x = (x_i)\) such that \(x_i \in K_i\) for all but a finite number of indices. \(G_0\) is said to be the local direct product of the \(G_i\) (relative to the \(K_i\) ); \(G_0 / K\) is then a discrete group isomorphic to the direct product of the \(G_i / K_i\) .

¶ 27) A group G is said to have the property P(n) (n an integer \textgreater0) if, for each system \((s_{1}, \ldots, s_{n})\) of n elements of G, there is a permutation \(\sigma \in S_{n}\) , other than the identity permutation (but depending on \(s_{1}, \ldots, s_{n}\) ) such that \(s_{\sigma(1)} \ldots s_{\sigma(n)} = s_{1} \ldots s_{n}\) .

\begin{enumerate}
\def\labelenumi{\alph{enumi})}
\item
  Show that the non-commutative group \(\mathfrak{S}_3\) satisfies \(\mathbf{P}(4)\) .
\item
  Let G be a connected Hausdorff topological group. Show that if G satisfies \(\mathrm{P}(n)\) for some integer n \textgreater{} 1, then G is commutative. {[}Show that G satisfies \(\mathrm{P}(n - 1)\) : given n - 1 elements \(s_{1}, \ldots, s_{n-1}\) of G, reduce to the case where there is a neighbourhood U of e such that, for each \(x \in U\) , we have \(s_{1} \ldots s_{n-1}x = s_{1} \ldots s_{i}xs_{i+1} \ldots s_{n-1}\) for some index \(i = i(x)\) ; deduce that the centralizer in G of one of the elements \(s_{j+1} \ldots s_{n-1} (0 \leqslant j \leqslant n - 2)\) is both open and closed and therefore equal to G. Hence the result if \(j \geqslant 1\) ; if j = 0, show that
\end{enumerate}

\[
\left. s _ {1} s _ {2} \dots s _ {n - 1} = s _ {2} \dots s _ {n - 1} s _ {1} \right].
\]

¶ 28) a) Let S be a stable subset of a topological group G. Show that \(\mathring{S}\) and \(\overline{S}\) are stable and that \(SS\subset\mathring{S}\) and \(\mathring{SS}\subset\mathring{S}\) . If also S is open and \(e\in\overline{S}\) , then \(S=\mathring{\overline{S}}\) (if \(x\in\overline{S}-S\) , show that \(x\) is a frontier point of S).

\begin{enumerate}
\def\labelenumi{\alph{enumi})}
\setcounter{enumi}{1}
\item
  From now on in this exercise, G is a commutative topological group, written additively. For every non-empty subset A of G, let \(s(A)\) denote the set of all \(x \in G\) such that \(x + A \subset A\) , i.e.~the intersection of the sets A - y as y runs through A; let \(b(A) = s(A) \cap s(-A)\) . The set \(s(A)\) is a stable subset of G, and \(b(A)\) is a subgroup of G consisting of all \(x \in G\) such that \(x + A = A\) . Show that if F is any non-empty closed subset of G, then \(s(F)\) is closed. For every non-empty subset A of G, we have \(s(A) \subset s(\overline{A})\) and \(s(A) \subset s(\mathring{A})\) ; and if \(A = \mathring{\overline{A}}\) , then \(s(\mathring{\overline{A}}) = s(\overline{\overline{A}})\) .
\item
  Let \(\mathfrak{S}\) be the set of all stable open subsets \(S\) of \(G\) such that \(0 \notin S\) , and suppose that \(\mathfrak{S} \neq \emptyset\) . Then the set \(\mathfrak{S}\) , ordered by inclusion, is inductive; let \(\mathfrak{M} \neq \emptyset\) be the set of its maximal elements. Show that if \(\mathbf{M} \in \mathfrak{M}\) , then \(\mathbf{M} = \overline{\mathbf{M}}\) {[}use \(a\) {]} and that, for each integer \(n > 0\) , \(\mathbf{M}\) is equal to the set of all \(x \in \mathbf{G}\) such that \(nx \in \mathbf{M}\) . In addition, the subgroup \(b(\mathbf{M})\) is equal to the complement of \(\mathbf{M} \cup (-\mathbf{M})\) in \(\mathbf{G}\) (if \(x \notin \mathbf{M}\) and \(-x \notin \mathbf{M}\) , consider the union of the sets \(kx + \mathbf{M}\) for integers \(k \geqslant 0\) ); deduce that \(\mathbf{G} = \mathbf{M} - \mathbf{M}\) .
\end{enumerate}

A closed subgroup \(\mathbf{B}\) of \(\mathbf{G}\) is said to be residual if there exists \(\mathbf{M} \in \mathfrak{M}\) such that \(\mathbf{B} = b(\mathbf{M})\) .

\begin{enumerate}
\def\labelenumi{\alph{enumi})}
\setcounter{enumi}{3}
\item
  The intersection of all the residual subgroups of G is called the radical of G and is denoted by \(T_{G}\) . If \(T_{G}=G\) (resp. \(T_{G}=\{0\}\) ) then G is said to be a radical group (resp. a group without radical). For x to belong to \(T_{G}\) it is necessary and sufficient that every stable open subset of G which contains x should also contain 0. Deduce that if G is discrete, then \(T_{G}\) is the torsion subgroup of G. If f is any continuous homomorphism of G into a commutative topological group \(G'\) , then \(f(T_{G})\subset T_{G'}\) . If H is a subgroup of G contained in \(T_{G}\) , then \(T_{G/H}=T_{G}/H\) . The subgroup \(T_{G}\) is the smallest of the subgroups H of G such that G/H is without radical.
\item
  Show that, among the subgroups H of G such that \(T_{H} = H\) , there is a largest, \(H_{0} \subset T_{G}\) , which is closed.
\end{enumerate}

\(f)\) If \(\mathbf{T}_{\mathbf{G}} \neq \mathbf{G}\) , then \(\mathbf{T}_{\mathbf{G}}\) is open if and only if there is a residual open subgroup in \(\mathbf{G}\) {[}show that if \(b(\mathbf{M})\) is open for some \(\mathbf{M} \in \mathfrak{M}\) , then \(\mathbf{T}_{\mathbf{G}} \supset b(\mathbf{M})\) , by using \(e)\) {]}.

* 29) Let G be the topological group R, let \(\varphi\) be the canonical mapping of R onto its quotient group \(\mathbf{T} = \mathbf{R} / \mathbf{Z}\) , and let \(\theta\) be an irrational number. The group G operates continuously on \(\mathbf{E} = \mathbf{T}^2\) by the rule \((s, (x, y)) \to (x + \varphi(s), y + \varphi(\theta s))\) . Show that the stabilizer of each \(z \in \mathbf{E}\) consists of 0 alone, that the orbit of \(z\) is dense in E and that it is not a topological homogeneous space relative to G.*

\begin{enumerate}
\def\labelenumi{\arabic{enumi})}
\setcounter{enumi}{29}
\tightlist
\item
  A topological group G is said to have no arbitrarily small subgroups if there is a neighbourhood V of e in G such that \(\{e\}\) is the only subgroup of G contained in V.
\end{enumerate}

\begin{enumerate}
\def\labelenumi{\alph{enumi})}
\item
  Let G be a topological group, H a normal subgroup of G. Show that if H and G/H have no arbitrarily small subgroups, then neither does G.
\item
  Deduce from a) that if \(H_{1}\) , \(H_{2}\) are two normal subgroups of a topological group G such that \(G/H_{1}\) and \(G/H_{2}\) have no arbitrarily small subgroups, then neither does \(\mathrm{G}/(\mathrm{H}_{1} \cap \mathrm{H}_{2})\) .
\end{enumerate}

¶ 31) Let G be a connected topological group, H a subgroup of G, U an open subset of G such that H = UH. Let \(V = H \cap (U^{-1}U)\); show that \(H = V^{\infty}\) {[}show that if \(A = H \cap C(V^{\infty})\), then the sets \(U.V^{\infty}\) and \(U.A\) do not intersect{]}.

